% Primeira subseção de Resultados, antes da análise dos sinais dos sensores.
\subsection{Caracterização da base de dados}
\label{subsec:caracterizacao_base_resultados}

A base analisada reúne 55.946 leituras em um arquivo CSV com 33 colunas.
Cada linha representa um registro temporal do nariz eletrônico durante
uma coleta de ar do espaço acima do solo de um vaso. Os registros incluem
os sinais dos seis sensores MQ, solo, temperatura, pressão, tempo,
identificação da coleta e classe experimental, além de colunas derivadas.
As leituras correspondem a 31 coletas realizadas em dez vasos, cinco do
grupo inoculado (classe 0) e cinco do grupo controle não inoculado
(classe 1). O número de leituras por coleta variou de 1.153 a 2.251,
com mediana de 2.005.

O grupo inoculado reuniu 26.342 leituras em 15 coletas, enquanto o grupo
controle reuniu 29.604 leituras em 16 coletas. A
Tabela~\ref{tab:distribuicao_base_resultados} apresenta a distribuição
das unidades experimentais e das leituras. Cada vaso foi medido em mais
de uma coleta; portanto, as 55.946 leituras não correspondem a 55.946
unidades biológicas independentes.

\begin{table}[htbp]
    \centering
    \caption{Distribuição das leituras e coletas na base analisada.}
    \label{tab:distribuicao_base_resultados}
    \begin{tabular}{lrrr}
        \hline
        Grupo & Vasos & Coletas & Leituras \\
        \hline
        Inoculado (classe 0) & 5 & 15 & 26.342 \\
        Controle (classe 1) & 5 & 16 & 29.604 \\
        \hline
        Total & 10 & 31 & 55.946 \\
        \hline
    \end{tabular}
    \vspace{0.1cm}
    {\footnotesize Fonte: Elaborado pelos autores (2026).}
\end{table}

Nos resultados de validação LOSO, o tamanho do treinamento variou de
48.612 a 52.412 leituras entre as dez rodadas; o conjunto de avaliação
de cada rodada continha de 3.534 a 7.334 leituras. Ao reunir as rodadas,
cada uma das 55.946 leituras recebeu uma previsão de avaliação.
Não houve um conjunto independente de validação nessa comparação.
